%% ============================================================
%%  §2 只选一层的代价：从 LLVM IR 到 XLA 与 Triton
%%  要证明的一句话：
%%    每个系统都把抽象层固定在某处；固定点越低则前端越厚、后端越薄，反之亦然。
%%    这是权衡曲线，不是优劣排序。
%%  素材：docs/05-section2-comparison.md
%% ============================================================

\section{只选一层的代价}
\label{sec:background}

\ref{sec:intro} 节把问题重述为"适配成本如何组织"。在讨论 MLIR 之前，需要先回答一个更基本的问题：
在 MLIR 出现之前，业界是如何选"那一层"的？本节考察五个代表性系统，只描述它们把抽象固定在哪里、
硬件知识从哪个口子进入，不做优劣判断——因为"选哪一层"本身是受使用场景约束的权衡，
而不是有唯一正确答案的设计题。

\subsection{五个系统，五种固定方式}

\paragraph{LLVM IR：把硬件知识整体推迟}
LLVM 的目标无关 SSA 中间表示在设计上就假定自己不含目标细节~\cite{llvm04}，
三地址码、标量或向量指令与显式内存语义构成其基本单位。硬件知识并非不存在，而是被
\textbf{整体推迟}到 IR 之后：指令选择、调度与 MC 层。因此扩展一个新后端，等于实现一套
\texttt{TargetLowering} 与指令选择模式；扩展一个新前端，则需要把源语言的全部语义下降成 LLVM IR。

\paragraph{Halide：把调度从算法里拆出来}
Halide 的核心贡献是算法与调度的解耦~\cite{halide12,halide13}：同一份算法可以配不同调度
（tile、vectorize、parallel），再配不同后端；官方支持 x86、ARM、MIPS、Hexagon、PowerPC、
RISC-V，以及 CUDA、OpenCL、OpenGL Compute、Metal、DirectX~12。这实际上是"把硬件知识外置成
调度"的早期形态，也是后来 \texttt{linalg} 的 structured operations 与 transform dialect 的思想来源。

\paragraph{TVM：固定两层}
TVM 同时暴露图级与算子级优化~\cite{tvm18}：图级负责算子融合等粗粒度变换，算子级负责
"映射到任意硬件原语"，并把内存延迟隐藏交给调度。硬件知识主要进入\textbf{算子级}——通过
tensor intrinsic、内存布局与 target 相关的下放规则。固定两层的好处是图级优化不必关心目标，
代价是接入新硬件往往要同时适配两层。

\paragraph{XLA：固定一层，其余交给后端}
XLA 的流程是三段式~\cite{xlaarch}：先在 StableHLO 上做\textbf{目标无关}的优化并转成内部
HLO；再交给后端做\textbf{目标相关}的 HLO 级优化；最后生成 LLVM IR 交给 LLVM 完成代码生成。
值得单独指出的是，布局（layout）在 HLO 内部就以显式 pass 的形式被指派~\cite{xlathunks}——
XLA 把"数据在内存里怎么摆"这件与硬件强相关的事，放在了与 IR 同一层处理。新增后端的路径是
"实现 \texttt{xla::Compiler} 接口"~\cite{xlabackend}。

XLA 官方文档自己给出了这个设计的动机，原话是：\emph{retargeting XLA should be significantly
simpler and more scalable than implementing every existing op from a frontend framework such as
PyTorch or TensorFlow for new hardware}~\cite{xlabackend}。这句话很重要：它说明"避免逐对适配"
不是本文事后附会的叙事，而是项目自己写下的设计目标。于是本文的问题不是新造的——同一问题上，
XLA 与 MLIR 给了两种不同的答案。

\paragraph{Triton：固定 tile 级，并且本身就是 MLIR 的一组方言}
Triton 面向分块（tiled）的神经网络计算~\cite{triton19}，把抽象固定在一块张量上，配合显式的
grid 模型。它常被误解为"绕开通用编译栈"的例子，但其编译器\textbf{本身就构建在 MLIR 上}：
其 IR 是一组方言，\texttt{tt}（Triton IR）$\rightarrow$ \texttt{ttg}（TritonGPU，含布局）
$\rightarrow$ \texttt{ttng}/\texttt{nvg}（NVIDIA）或 \texttt{amdg}（AMD）~\cite{tritondialects}。
换句话说，Triton 是 MLIR 的一种\textbf{纵向用法}——只保留必要的一两级方言，把 tile 级钉死，
再一路下到厂商方言；它与 IREE、StableHLO 这类"宽栈"用法共同构成 MLIR 生态的两种形态。

\begin{table}[t]
  \caption{各系统把抽象层固定在何处，以及硬件知识与新前端的进入点。}
  \label{tab:layer-fixing}
  \small
  \begin{tabularx}{\linewidth}{@{}l
      >{\raggedright\arraybackslash}p{1.5cm}
      >{\raggedright\arraybackslash}p{1.9cm}
      >{\raggedright\arraybackslash}p{1.7cm}
      >{\raggedright\arraybackslash}p{1.9cm}
      >{\raggedright\arraybackslash}p{1.3cm}@{}}
    \toprule
    系统 & 固定的抽象层 & 硬件知识进入点 & 新增前端 & 新增后端 & 复用单位 \\
    \midrule
    LLVM IR & 一层：目标无关 SSA IR & 指令选择、调度、MC 层
      & 生成 LLVM IR & 实现 TargetLowering / ISel & pass \\
    Halide & 一层：算法 + 调度 & 调度原语与后端
      & 写算法（内嵌 DSL） & 接入一个后端 & 调度原语 \\
    TVM & 两层：图级 + 算子级 & 算子级调度与 target 下放
      & 导入为图级 IR & target 描述 + codegen / BYOC & 调度原语 \\
    XLA & 一层：HLO & HLO 内部 pass 与后端重写
      & 生成 StableHLO & 实现 \texttt{Compiler} 接口 & HLO pass \\
    Triton & 一层：tile 级 & tile 布局（\texttt{ttg}）与厂商方言
      & 生成 tile 级程序 & 厂商方言 + lowering & dialect \\
    MLIR & \textbf{不固定} & 任意 dialect 边界
      & 定义前端 dialect & 定义目标 dialect + 转换 & dialect / 转换 \\
    \bottomrule
  \end{tabularx}
\end{table}

\subsection{三点评注}

\paragraph{一：层的位置是守恒的，不是可以省掉的}
表中六行没有一行是"无层"的——每行都选了至少一层，TVM 选了两层。能变的只是"选在哪"。
因此严格的说法不是"有了 MLIR 就不必选层"，而是：\textbf{抽象层的位置可以随域调整，
但抽象层的总量不会消失}。MLIR 改变的是"谁来做这个选择、以及能否在同一程序里保留多个选择"，
而不是"可以不做选择"。

\paragraph{二：固定点越低，前端越厚、后端越薄——这是一条权衡曲线}
LLVM IR 固定在最低处：前端必须承担全部语义下降，后端只需做指令选择；XLA 固定在高处：
前端只需产出 StableHLO（极薄），后端必须承担全部目标相关重写（极厚）。Halide、TVM、Triton
落在曲线的不同高度上。

这条曲线与 \ref{sec:intro} 的成本模型直接咬合：\textbf{最优固定点由 $N$ 与 $M$ 的相对大小决定}。
若一个系统预期被很多前端复用、只服务少数后端，把层压低更划算；反之应把层抬高。
这个判断与具体技术优劣无关，只与生态规模有关。

\paragraph{三：XLA 与 MLIR 是同一问题的两种答案，差别在"谁承担不确定性"}
两者都以"避免为每个硬件重写全部算子"为动机（前者见上引 XLA 原文，后者见 MLIR 的
\emph{software fragmentation}~\cite{mlir21}）。差别在于不确定性被放在哪里：XLA 把它收进后端，
以换取 HLO 这一层的稳定；MLIR 把它放到 dialect 边界上，允许边界本身被重新协商。

这一区别不需要评判谁更好，但它解释了第~\ref{sec:evidence} 节的实证结果：如果 MLIR 的成功
取决于"边界能否被协商成功"，那么\textbf{边界协商的成本}就是关键变量——这正是第~\ref{sec:contract} 节
要考察的内容。

\paragraph{关于表~\ref{tab:layer-fixing} 的 MLIR 一行}
该行是\textbf{机制描述}（依据 MLIR 的设计文献与官方文档），不是本文的实测结果；
它与第~\ref{sec:evidence} 节基于语料统计的证据强度不同，不应等同看待。
